\begin{helper}
Let \(I\subseteq \operatorname{Fin}(n)\), let \(H\) be a history cell, let
\(T\) be a target event, let \(U\) be the terminal-coordinate set, and let
\(\omega_0\) be the representative sample for the history cell.  Assume:
\[
|I|=k,
\]
every \(\omega\in H\) has the same embedding as \(\omega_0\), every
\(\omega\in H\) satisfies the projection-size event with parameters
\((k,\ell_R,\ell_B)\), \(T\subseteq H\), every staged-open coordinate of a
sample in \(H\) lies in \(U\), and for every \(\omega\in T\) the blue present
staged-open support
\[
B(\omega)=
\{e\in \operatorname{StagedOpen}(\omega,\varepsilon,I):
  X_e(\omega)=1\text{ and }e\text{ is a blue coordinate}\}
\]
has cardinality \(u_B\).  Then there are a natural number \(N_B\) and a
finite family \(\mathcal B\) of terminal blue support labels such that
\[
N_B\le k^2,\qquad |\mathcal B|\le \binom{N_B}{u_B}.
\]
For every \(B\in\mathcal B\), \(B\) has cardinality \(u_B\), satisfies
\(B\subseteq U\), and consists only of blue coordinates.  For every
\(\omega\in T\), the set \(B(\omega)\) belongs to \(\mathcal B\).
\end{helper}
\begin{proof}
The proof constructs \(\mathcal B\) from a fixed blue projection universe.
Because the history cell fixes the embedding, the blue projection of each
vertex of \(I\) is the same in every sample of \(H\) as in the representative
sample \(\omega_0\).  Let \(I_B\) be this fixed blue projection of \(I\).
The projection-size hypothesis is the local tablet hypothesis which makes
this projected universe available uniformly over the cell.

Let \(U_B\) be the set of blue terminal coordinates obtained from projected
pairs of vertices in \(I_B\).  Since \(|I|=k\), the crude ordered-pair bound
gives
\[
|U_B|\le |I_B|^2\le k^2.
\]
Set \(N_B=|U_B|\).  Define \(\mathcal B\) to be the finite family of all
\(u_B\)-element subsets of \(U_B\) which are also contained in \(U\).  This is
a subfamily of the \(u_B\)-subsets of an \(N_B\)-element set, hence
\[
|\mathcal B|\le \binom{N_B}{u_B}.
\]
Every member of \(\mathcal B\) has cardinality \(u_B\), lies in \(U\), and is
blue because \(U_B\) contains only blue coordinates.

Now fix \(\omega\in T\).  Since \(T\subseteq H\), every staged-open coordinate
of \(\omega\) lies in \(U\).  The set \(B(\omega)\) is exactly the filter of
the staged-open set by present blue coordinates, so it is a subset of \(U\).
The fixed-embedding argument places every such blue coordinate in \(U_B\), and
the fixed-count hypothesis gives \(|B(\omega)|=u_B\).  Therefore
\(B(\omega)\in\mathcal B\), proving the label-covering assertion.
\end{proof}
